\chapter{Action Representation Is the Control Boundary}

\section*{학습 목표}

이 장은 책 전체의 action API를 고정한다. VLA에서 가장 중요한 질문은 어떤 backbone을 쓰는가만이 아니다. 더 근본적인 질문은 action을 무엇으로 표현하는가이다. 로봇 action은 token일 수도 있고, continuous vector일 수도 있고, waypoint, delta pose, action chunk, diffusion/flow distribution, whole-body command일 수도 있다. 이 선택은 출력 포맷의 문제가 아니라 VLA policy layer와 motion/control layer 사이의 책임 분할을 정하는 control boundary 문제다.

\begin{enumerate}
  \item action representation을 control boundary로 해석한다.
  \item token, vector, waypoint, chunk, distribution, action expert를 구분한다.
  \item 각 표현이 latency, safety, recovery, controller compatibility에 미치는 영향을 설명한다.
  \item 실제 robot task에 맞는 action interface를 고르는 checklist를 만든다.
\end{enumerate}

\section{Why action representation is the control boundary}

VLA policy를 가장 추상적으로 쓰면 다음과 같다.
\begin{equation}
  \act_t
  \sim
  \pi_{\theta}
  (\cdot \mid \obs_{\leq t},\ell,\conf_t).
  \label{eq:ch13-single-action-policy}
\end{equation}
여기서 $\act_t$가 무엇인지가 이 장의 핵심이다. $\act_t$가 discrete token이면 decoder와 tokenizer가 중요하다. $\act_t$가 end-effector delta pose이면 frame convention, Cartesian servo, IK가 중요하다. $\act_t$가 action chunk이면 horizon $H$, update frequency, interruption rule이 중요하다. $\act_t$가 torque라면 stability와 safety assurance 부담이 급격히 커진다.

따라서 action representation은 다음 계약을 정의한다.
\begin{equation}
  \mathrm{Contract}
  =
  (
  \mathcal{A},
  f_{VLA},
  f_{ctrl},
  \mathcal{C}_{safe},
  \mathcal{E}_{exec}
  ),
  \label{eq:ch13-action-contract}
\end{equation}
여기서 $\mathcal{A}$는 action space, $f_{VLA}$는 VLA update rate, $f_{ctrl}$은 controller loop rate, $\mathcal{C}_{safe}$는 safety constraint, $\mathcal{E}_{exec}$는 executor/controller이다. 같은 성능표를 가진 두 VLA라도 이 contract가 다르면 실제 시스템은 전혀 다르다.

\begin{tcolorbox}[title={Roboticist's takeaway}]
Action representation은 VLA의 출력 포맷이 아니라, 의미 기반 policy와 물리 controller 사이의 계약이다.
\end{tcolorbox}

\section{Action representation spectrum}

이 책에서는 다음 spectrum을 기본 taxonomy로 사용한다.
\begin{equation}
  \begin{aligned}
  \mathcal{A}_{VLA}
  =
  \{&
  \mathrm{token},
  \mathrm{vector},
  \mathrm{waypoint},
  \mathrm{delta\ pose},\\
  &
  \mathrm{chunk},
  \mathrm{distribution},
  \mathrm{expert},
  \mathrm{whole\ body}
  \}.
  \end{aligned}
  \label{eq:ch13-action-taxonomy}
\end{equation}

\begin{table}[h]
  \centering
  \caption{Action representation taxonomy}
  \begin{tabularx}{\linewidth}{p{0.18\linewidth}X X}
    \toprule
    Type & Downstream interface & Main failure mode \\
    \midrule
    Action token & token decoder, post-decoder, action wrapper & quantization, invalid token, decoding latency \\
    Continuous vector & servo, OSC, impedance, action wrapper & scaling error, saturation, mode averaging \\
    Waypoint & IK, motion planner, MPC & waypoint가 trajectory feasibility를 보장하지 않음 \\
    Delta pose & Cartesian servo, OSC, impedance & frame convention error, drift, limit violation \\
    Action chunk & chunk executor, low-level tracker & chunk inertia, interruption difficulty \\
    Distribution & sampling, diffusion/flow action head & sampling cost, safety projection difficulty \\
    Action expert & VLM backbone plus motor head & semantic-motor mismatch, fine-tuning burden \\
    Whole-body action & WBC, balance, locomotion controller & balance, self-collision, contact instability \\
    \bottomrule
  \end{tabularx}
  \label{tab:ch13-taxonomy}
\end{table}

이 spectrum은 순위가 아니다. 왼쪽으로 갈수록 semantic abstraction이 크고, 오른쪽으로 갈수록 physical responsibility가 커진다. 좋은 interface는 task, robot, controller, latency, data, safety envelope에 따라 달라진다.

\section{Canonical VLA case comparison}

최근 VLA 흐름을 action representation 관점에서 다시 읽으면, 모델 family의 차이는 backbone 크기보다 $\act_t$가 어떤 형식으로 생성되고 어떤 controller boundary를 통과하는가에서 선명해진다. RT-1과 RT-2는 action token을 통해 language model training recipe와 robot action을 연결했다 \cite{rt12022,rt22023}. OpenVLA는 open VLA baseline을 제공했고, OpenVLA-OFT는 continuous action, action chunking, parallel decoding, fine-tuning objective가 inference speed와 success를 바꾼다는 점을 보여준다 \cite{openvla2024,openvlaoft2025}. $\pi_0$는 flow matching action expert를 사용하고, FAST는 action tokenizer 자체를 고주파 robot action sequence의 압축 문제로 다룬다 \cite{pi02024,fast2025}. RDT-1B 같은 diffusion transformer 계열은 bimanual manipulation에서 unified action space와 diffusion-based action generation을 결합하는 또 다른 action-interface 사례로 읽을 수 있다 \cite{rdt1b2024}.

\begin{table}[h]
  \centering
  \footnotesize
  \setlength{\tabcolsep}{3.5pt}
  \renewcommand{\arraystretch}{1.12}
  \caption{Action representation 관점에서 본 대표 VLA 계열}
  \begin{tabularx}{\linewidth}{@{}>{\RaggedRight\arraybackslash}p{0.16\linewidth}>{\RaggedRight\arraybackslash}p{0.25\linewidth}>{\RaggedRight\arraybackslash}p{0.24\linewidth}>{\RaggedRight\arraybackslash}X@{}}
    \toprule
    Model family & Action representation & Controller interface & Latency / safety note \\
    \midrule
    RT-1 & discrete action token & robot-specific wrapper & token rate, decoding cost, quantization, invalid action mapping \\
    RT-2 & action as token & action tokenizer/wrapper & autoregressive delay, precision, smoothness, vocabulary restriction \\
    OpenVLA & token/action decoding & fine-tuned robot policy & large-model throughput, target-robot adaptation, calibration \\
    OpenVLA-OFT & continuous action + chunking & faster decoder/wrapper & fast chunk generation, interruption rule, stale action validation \\
    $\pi_0$ & flow matching action expert & high-frequency action generation & denoising budget, control budget, action distribution validation \\
    FAST & compressed action tokenizer & token-to-action decoder & efficient AR generation, reconstruction error, control fidelity \\
    \bottomrule
  \end{tabularx}
  \label{tab:ch13-vla-action-cases}
\end{table}

표~\ref{tab:ch13-vla-action-cases}의 목적은 model ranking이 아니다. 같은 VLA라도 action이 token인지, continuous vector인지, chunk인지, flow-generated trajectory인지에 따라 downstream controller, safety shield, evaluation metric이 달라진다. 그러므로 VLA 논문을 읽을 때 action representation table이 없으면 그 논문은 아직 robot system claim을 충분히 닫지 않은 것이다.

\begin{casebox}{Case study: OpenVLA-OFT, FAST, $\pi_0$를 action boundary로 읽기}
OpenVLA-OFT, FAST, $\pi_0$는 모두 ``VLA 성능 개선''으로 묶일 수 있지만, Real VLA 관점에서는 서로 다른 병목을 겨냥한다 \cite{openvlaoft2025,fast2025,pi02024}. OpenVLA-OFT는 continuous action, chunking, faster decoding, fine-tuning objective를 통해 target robot adaptation과 speed/success trade-off를 다룬다. FAST는 high-frequency robot action sequence를 token compression 문제로 읽는다. $\pi_0$는 flow matching action expert를 사용해 action distribution 자체를 생성한다.

\begin{center}
\small
\begin{tabularx}{\linewidth}{@{}p{0.22\linewidth}X X@{}}
\toprule
사례 & 실제로 바뀐 것 & 반드시 물어야 할 질문 \\
\midrule
OpenVLA-OFT & action decoding, continuous action, chunking & chunk interruption과 stale action rejection은 어떻게 되는가? \\
FAST & action tokenizer와 reconstruction fidelity & token compression error가 controller tracking error로 어떻게 나타나는가? \\
$\pi_0$ & flow-based action generation & denoising budget과 robot control budget이 동시에 만족되는가? \\
\bottomrule
\end{tabularx}
\end{center}

이 세 사례의 공통점은 backbone 논쟁보다 action API 논쟁이 더 중요하다는 점이다. 논문을 읽을 때 ``무슨 모델인가''보다 ``어떤 action을 어느 rate로 내고, 어떤 controller와 safety layer가 그것을 받아들이는가''를 먼저 확인해야 한다.
\end{casebox}

\section{Discrete action tokens}

Discrete action token은 robot action을 vocabulary item처럼 다룬다. 연속 action $\act_t\in\mathbb{R}^d$를 binning 또는 tokenizer로 변환하면
\begin{equation}
  z_t
  =
  \mathrm{Tok}(\act_t),
  \qquad
  z_t\in\mathcal{V}_{act}.
  \label{eq:ch13-tokenizer}
\end{equation}
VLA는 다음 token을 예측한다.
\begin{equation}
  p_{\theta}(z_t \mid \obs_{\leq t},\ell,z_{<t}).
  \label{eq:ch13-action-token-prob}
\end{equation}
이 방식은 VLM/LLM training pipeline과 잘 맞는다. cross-entropy loss, autoregressive decoding, text token과 action token의 통합이 가능하기 때문이다.

하지만 token action에는 quantization error가 있다.
\begin{equation}
  \epsilon_{q}
  =
  \|\act_t - \mathrm{Dec}(\mathrm{Tok}(\act_t))\|.
  \label{eq:ch13-quantization-error}
\end{equation}
fine manipulation, dexterous hand, force-sensitive contact에서는 작은 quantization error가 task failure로 이어질 수 있다. 또한 action token sequence가 길어지면 decoding latency가 증가한다. 따라서 token action은 semantic transfer와 training convenience가 강한 대신, high-frequency continuous control에서는 controller나 post-decoder의 보완이 필요하다.

\section{Continuous regression and normalization}

Continuous action vector는 VLA가 실수 벡터를 직접 예측하는 방식이다.
\begin{equation}
  \hat{\act}_t
  =
  f_{\theta}(\obs_{\leq t},\ell,\conf_t),
  \qquad
  \hat{\act}_t\in\mathbb{R}^{d}.
  \label{eq:ch13-continuous-regression}
\end{equation}
behavior cloning에서는 주로 다음 손실을 쓴다.
\begin{equation}
  \mathcal{L}_{reg}
  =
  \sum_t
  \|\hat{\act}_t - \act_t^{E}\|_1
  \quad
  \mathrm{or}
  \quad
  \sum_t
  \|\hat{\act}_t - \act_t^{E}\|_2^2.
  \label{eq:ch13-regression-loss}
\end{equation}

continuous action은 controller와 자연스럽게 연결된다. 예를 들어 normalized end-effector delta command를 Cartesian servo에 넣을 수 있다. 그러나 normalization이 잘못되면 같은 vector가 robot마다 전혀 다른 물리 의미를 갖는다.
\begin{equation}
  \act_t^{phys}
  =
  \sigma_a \odot \act_t^{norm}
  +
  \mu_a.
  \label{eq:ch13-action-denormalization}
\end{equation}
$\mu_a,\sigma_a$는 dataset, robot, controller mode에 의존한다. multi-embodiment dataset에서는 이 값이 성능과 안전성에 직접 영향을 준다.

continuous regression의 또 다른 문제는 mode averaging이다. 가능한 action mode가 여러 개인데 평균을 내면 실제로는 어떤 mode도 아닌 action이 나온다.
\begin{equation}
  \hat{\act}
  =
  \mathbb{E}[\act\mid\obs,\ell]
  \notin
  \{\act^{(1)},\act^{(2)},\ldots\}.
  \label{eq:ch13-mode-averaging}
\end{equation}
contact-rich manipulation에서는 mode averaging이 slip, collision, failed insertion으로 이어질 수 있다. 이것이 diffusion/flow action distribution이 중요해지는 이유다.

\section{Waypoint, delta pose, and trajectory target}

Waypoint와 delta pose는 continuous vector의 특수한 interface이다. end-effector pose를 $T_t\in SE(3)$라고 하자. delta pose command는 다음 target을 만든다.
\begin{equation}
  T_{t+1}^{target}
  =
  T_t \exp(\Delta \xi_t),
  \label{eq:ch13-delta-pose}
\end{equation}
여기서 $\Delta \xi_t$는 twist coordinate로 표현된 상대 motion이다. 이 command는 Cartesian servo, operational-space control, impedance controller와 연결되기 쉽다.

Waypoint는 더 큰 시간 단위의 reference이다.
\begin{equation}
  r_t
  =
  (p_t^{target},R_t^{target},g_t^{gripper}),
  \label{eq:ch13-waypoint}
\end{equation}
하지만 waypoint는 trajectory가 아니다. waypoint가 collision-free path, joint-limit satisfaction, dynamic feasibility를 보장하지 않는다. 따라서 waypoint interface는 motion planner 또는 MPC와 함께 써야 한다.

Trajectory target은 reference sequence이다.
\begin{equation}
  r_{t:t+H}
  =
  (r_t,r_{t+1},\ldots,r_{t+H}).
  \label{eq:ch13-trajectory-target}
\end{equation}
이때 VLA는 full low-level command를 내는 것이 아니라 controller가 추적할 목표를 낸다. 이 차이를 명확히 하지 않으면 논문에서 말하는 action과 실제 robot driver command가 혼동된다.

\section{Action chunking as temporal interface}

Action chunk는 한 시점 action이 아니라 여러 time step의 action sequence를 한 번에 예측하는 출력 단위이다.
\begin{equation}
  A_t^{(H)}
  =
  (\act_t,\act_{t+1},\ldots,\act_{t+H-1})
  \sim
  \pi_{\theta}
  (\cdot \mid \obs_{\leq t},\ell).
  \label{eq:ch13-action-chunk}
\end{equation}
chunking은 단순히 inference를 덜 자주 하기 위한 trick이 아니다. temporal abstraction이다. policy가 짧은 motion segment를 한 번에 만들면 single-step compounding error를 줄이고, bimanual manipulation처럼 coordination이 필요한 task에서 temporal consistency를 얻을 수 있다.

하지만 chunk horizon $H$는 trade-off를 만든다. $H$가 길수록 motion은 부드러워질 수 있지만, unexpected contact나 perception update가 들어왔을 때 빠르게 수정하기 어렵다. 이를 chunk inertia라고 부를 수 있다.
\begin{equation}
  \mathrm{RecoveryDelay}
  \approx
  \frac{H_{remain}}{f_{ctrl}}
  +
  \frac{1}{f_{VLA}}.
  \label{eq:ch13-recovery-delay}
\end{equation}
따라서 chunk executor는 interruption rule을 가져야 한다.
\begin{equation}
  \mathrm{Exec}(A_t^{(H)})
  =
  \begin{cases}
  \mathrm{continue}, & \mathrm{safe}(x_\tau,A_t^{(H)})=1,\\
  \mathrm{interrupt}, & \mathrm{safe}(x_\tau,A_t^{(H)})=0.
  \end{cases}
  \label{eq:ch13-chunk-interrupt}
\end{equation}

\section{Diffusion policies and multimodal actions}

Continuous regression은 여러 가능한 action mode를 평균낼 수 있다. Diffusion policy는 action chunk distribution을 denoising process로 표현한다. 단순화하면 noisy action $A^K$에서 시작해 denoising을 반복한다.
\begin{equation}
  A^{k-1}
  =
  D_{\theta}(A^k,k,\obs,\ell),
  \qquad
  k=K,K-1,\ldots,1.
  \label{eq:ch13-diffusion-denoise}
\end{equation}
최종 action chunk는 $A^0$이다.
\begin{equation}
  A_t^{(H)}
  =
  A^0.
  \label{eq:ch13-diffusion-output}
\end{equation}

이 방식의 장점은 multimodal action distribution을 표현하기 쉽다는 점이다. 예를 들어 object를 왼쪽에서 잡을 수도 있고 오른쪽에서 잡을 수도 있다. 두 mode의 평균은 실패할 수 있지만, distributional generator는 하나의 mode를 sample할 수 있다.

그러나 diffusion은 sampling cost를 갖는다.
\begin{equation}
  T_{infer}
  \approx
  K \cdot T_{denoise}.
  \label{eq:ch13-diffusion-latency}
\end{equation}
real-time robot에서는 $T_{infer}$가 controller loop와 맞아야 한다. 또한 sampling diversity는 safety guarantee가 아니다. sample된 action은 여전히 collision, joint limit, force limit, workspace constraint를 통과해야 한다.

\section{Flow matching and action expert}

Flow matching은 action space에서 continuous transformation을 학습하는 방식으로 볼 수 있다. noise 또는 prior sample $A_0$에서 data action $A_1$으로 흐르는 vector field를 학습한다.
\begin{equation}
  \frac{dA_{\tau}}{d\tau}
  =
  v_{\theta}(A_{\tau},\tau,\obs,\ell),
  \qquad
  \tau\in[0,1].
  \label{eq:ch13-flow-matching}
\end{equation}
이 구조는 VLM backbone과 continuous action expert를 분리하기 좋다.
\begin{equation}
  h^{vlm}
  =
  F_{\theta}(\obs,\ell),
  \qquad
  A_t^{(H)}
  =
  G_{\phi}(h^{vlm},\obs,\conf_t).
  \label{eq:ch13-action-expert}
\end{equation}
$F_{\theta}$는 semantic representation을 만들고, $G_{\phi}$는 motor action distribution을 만든다. 이 분리는 Real VLA 관점에서 중요하다. semantic backbone과 motor head의 학습 속도, data requirement, deployment constraint가 다르기 때문이다.

하지만 action expert는 controller와 맞아야 한다. expert가 만든 chunk가 downstream controller의 rate, saturation, smoothness, frame convention과 맞지 않으면 좋은 semantic representation도 실제 robot에서는 실패한다.

\section{Action tokenizers and compression}

Action tokenization은 단순 binning으로 끝나지 않는다. robot action은 시간 구조와 주파수 구조를 갖는다. action trajectory $A_t^{(H)}$를 compact token $z_t$로 압축한다고 하자.
\begin{equation}
  z_t
  =
  \mathrm{Tok}(A_t^{(H)}),
  \qquad
  \hat{A}_t^{(H)}
  =
  \mathrm{Dec}(z_t).
  \label{eq:ch13-chunk-tokenization}
\end{equation}
좋은 tokenizer는 reconstruction error뿐 아니라 smoothness와 control relevance를 보존해야 한다.
\begin{equation}
  \mathcal{L}_{tok}
  =
  \|A-\hat{A}\|^2
  +
  \lambda_s
  \|\nabla A-\nabla \hat{A}\|^2
  +
  \lambda_f
  \|\Phi(A)-\Phi(\hat{A})\|^2.
  \label{eq:ch13-tokenizer-loss}
\end{equation}
여기서 $\Phi$는 frequency-domain feature 또는 task-relevant feature이다. high-frequency dexterous behavior를 모두 제거한 token은 language model에는 편해도 robot control에는 부적합할 수 있다.

\section{Choosing an interface}

Action representation 선택은 다음 조건으로 결정해야 한다.
\begin{table}[h]
  \centering
  \caption{Task별 action interface 선택}
  \begin{tabularx}{\linewidth}{p{0.23\linewidth}X X}
    \toprule
    Task / embodiment & Preferred representation & Safety/evaluation focus \\
    \midrule
    Tabletop pick-and-place & delta pose, waypoint, short chunk & collision, grasp failure, pose tracking \\
    Bimanual manipulation & action chunk, continuous vector, distribution & coordination, chunk interruption, contact force \\
    Dexterous manipulation & diffusion/flow distribution, tactile-conditioned chunk & slip, multimodality, latency, force limit \\
    Mobile manipulation & subgoal plus waypoint plus local chunk & base-arm coordination, obstacle proximity \\
    Humanoid whole-body task & whole-body reference, hierarchy, WBC target & balance, self-collision, fall recovery \\
    Human-proximity task & conservative waypoint or skill call with strong shield & intervention, E-stop, audit log \\
    \bottomrule
  \end{tabularx}
  \label{tab:ch13-design-matrix}
\end{table}

이 표의 핵심은 universal best action representation은 없다는 것이다. token이 항상 구식인 것도 아니고, diffusion이나 flow가 항상 정답인 것도 아니다. task와 controller가 요구하는 물리적 계약에 맞아야 한다.

\section{Failure modes}

\begin{table}[h]
  \centering
  \caption{Action representation failure modes}
  \begin{tabularx}{\linewidth}{p{0.25\linewidth}X X}
    \toprule
    Failure mode & Mechanism & Mitigation \\
    \midrule
    Quantization loss & discrete bin이 fine action을 파괴 & finer tokenizer, residual controller, continuous head \\
    Mode averaging & regression이 여러 mode의 평균을 냄 & diffusion/flow, mixture model, goal-conditioned disambiguation \\
    Chunk inertia & 긴 chunk가 빠른 recovery를 막음 & interruption rule, shorter horizon, monitor-triggered replanning \\
    Controller mismatch & action type과 downstream controller가 맞지 않음 & explicit action schema, controller wrapper, rate matching \\
    Latency mismatch & $f_{VLA}$가 $f_{ctrl}$보다 너무 낮음 & asynchronous execution, chunking, fallback controller \\
    Embodiment mismatch & 같은 vector가 robot마다 다른 물리 의미를 가짐 & embodiment adapter, normalization card, action calibration \\
    Safety incompatibility & shield가 action을 해석하지 못함 & safety-readable action type, projection, runtime monitor \\
    \bottomrule
  \end{tabularx}
  \label{tab:ch13-failure-modes}
\end{table}

\section{Checklist}

새 VLA system을 읽거나 설계할 때는 다음 질문에 답해야 한다.
\begin{enumerate}
  \item action은 token, vector, waypoint, delta pose, chunk, distribution, torque 중 무엇인가?
  \item action은 어느 frame에서 정의되는가?
  \item VLA update rate $f_{VLA}$와 controller rate $f_{ctrl}$은 얼마인가?
  \item action normalization과 clipping rule은 무엇인가?
  \item safety shield는 action 전, action 후, controller reference 후 중 어디에서 작동하는가?
  \item action chunk가 있다면 horizon $H$와 interruption rule은 무엇인가?
  \item diffusion/flow를 쓴다면 sampling time과 safety projection time은 얼마인가?
  \item multi-embodiment라면 action adapter 또는 embodiment descriptor가 있는가?
  \item evaluation은 success rate뿐 아니라 latency, smoothness, saturation, recovery를 측정하는가?
\end{enumerate}

\section{Running example: mug, rack, and glass}

``머그컵을 drying rack에 올리고 유리컵은 건드리지 말라''는 같은 instruction이라도 action representation에 따라 전혀 다른 system claim이 된다. RT-style token action은 ``move-left'', ``close-gripper'' 같은 discrete command로 표현될 수 있지만, rack slot과 유리컵 사이 간격이 작으면 quantization error가 collision margin을 먹을 수 있다. Continuous delta pose는 제어기와 잘 연결되지만, grasp approach가 여러 mode를 가질 때 평균 action이 rack edge를 향할 수 있다. Action chunk는 pick--lift--transport 구간을 부드럽게 만들 수 있지만, 유리컵이 흔들리거나 사람이 끼어들면 interruption rule이 없을 때 위험하다. Diffusion 또는 flow action expert는 여러 가능한 approach trajectory를 표현할 수 있지만, sampling된 chunk가 forbidden region을 침범하지 않는지 safety shield가 해석할 수 있어야 한다.

\begin{definitionbox}{이 예제에서 action representation이 남겨야 할 기록}
좋은 VLA report는 mug pose uncertainty, glass forbidden region, rack placement target, chosen action type, controller reference, chunk horizon, interruption rule, latency, safety rejection 횟수를 함께 남겨야 한다. ``성공했다''는 말만으로는 action representation의 공학적 품질을 평가할 수 없다.
\end{definitionbox}

\chapterwork
{RT-1/RT-2, OpenVLA-OFT, FAST, $\pi_0$, RDT-1B를 action representation contract로 비교하라 \cite{rt12022,rt22023,openvlaoft2025,fast2025,pi02024,rdt1b2024}.}
{action representation이 output format이 아니라 control boundary라는 말은 무엇을 의미하는가?}
{하나의 contact-rich task에 대해 token, continuous delta, waypoint, chunk, diffusion/flow action의 장단점을 controller/safety/evaluation 관점에서 비교하라.}

\section{요약}

Action representation은 Real VLA의 중심 문제다. token, continuous vector, waypoint, delta pose, chunk, diffusion/flow action은 각각 controller interface, latency, recovery granularity, safety projection 조건을 바꾼다. 따라서 좋은 VLA는 좋은 backbone만으로 정의되지 않는다. 좋은 VLA는 자신이 어떤 action contract 위에서 동작하는지 명시한다. 다음 장에서는 이 action representation이 어떤 data engine에서 학습되고 유지되는지 다룬다.
